% Part: applied-modal-logics
% Chapter: epistemic-logic
% Section: language-epistemic-logic

\documentclass[../../../include/open-logic-section]{subfiles}

\begin{document}

\olfileid{aml}{el}{pal}

\olsection{బహిరంగ ప్రకటన తర్కం}

కర్తల జ్ఞానం కాలక్రమంలో లేదా వారికి కొత్త సమాచారం
అందినప్పుడు ఎలా మారుతుందో గతి జ్ఞానసంబంధ తర్కాలు
వ్యక్తీకరించగలవు. వీటిలో అనేకం సమాచార
\emph{ఘటనలు} లేదా \emph{నవీకరణల} ద్వారా జ్ఞాన
మార్పులను చూపుతాయి. అత్యంత ప్రాథమిక నవీకరణ
బహిరంగ ప్రకటన: ఏదో సూత్రాన్ని సత్యంగా ప్రకటించడం,
ఆ ప్రకటన జరగడాన్ని కర్తలందరూ కలిసి గమనించడం.
దీన్ని చూపడానికి భాషను ఇలా విస్తరిస్తాం.

\begin{defn}
$G$ కర్త-సంకేతాల సమితి అనుకుందాం. బహిరంగ
ప్రకటనలతో కూడిన బహు-కర్త జ్ఞానసంబంధ తర్కపు
ప్రాథమిక భాషలో ఇవి ఉంటాయి:
\begin{enumerate}
  \tagitem{prvFalse}{\tetoken{అసత్యం}{!!{falsity}} కోసం ప్రతిజ్ఞావాక్య స్థిరాంకం~$\lfalse$.}{}
  \tagitem{prvTrue}{\tetoken{సత్యం}{!!{truth}} కోసం ప్రతిజ్ఞావాక్య స్థిరాంకం~$\ltrue$.}{}
  \item \tetoken{ప్రతిజ్ఞావాక్య చరాల}{!!{propositional variable}s}
    \tetoken{అనంతంగా లెక్కించదగిన}{!!{denumerable}s} సమితి: $\Obj
    p_0$, $\Obj p_1$, $\Obj p_2$, \dots
  \item ప్రతిజ్ఞావాక్య సంధాయకాలు: \startycommalist
  \iftag{prvNot}{\ycomma $\lnot$ (నిరాకరణ)}{}%
  \iftag{prvAnd}{\ycomma $\land$ (సంయోజనం)}{}%
  \iftag{prvOr}{\ycomma $\lor$ (విడియోజనం)}{}%
  \iftag{prvIf}{\ycomma $\lif$ (\tetoken{సోపాధికం}{!!{conditional}})}{}%
  \item $a \in G$ అయినప్పుడు జ్ఞాన కారకం $\Knows_a$.
  \item $!B$ \tetoken{ఒక సూత్రం}{!!a{formula}} అయినప్పుడు బహిరంగ ప్రకటన కారకం $[!B]$.
\end{enumerate}
\end{defn}

బహిరంగ ప్రకటన కారకం బాక్స్ కారకంలా పనిచేస్తుంది;
దానికి అనుగుణంగానే భాషకు ఆగమన నిర్వచనం ఇస్తాం:

\begin{defn}
  జ్ఞానసంబంధ భాషలోని \emph{\tetoken{సూత్రాలను}{!!^{formula}s}}
  ఇలా ఆగమన పద్ధతిలో నిర్వచిస్తాం:
  \begin{enumerate}
  \tagitem{prvFalse}{$\lfalse$ ఒక అణు \tetoken{సూత్రం}{!!{formula}}.}{}

  \tagitem{prvTrue}{$\ltrue$ ఒక అణు \tetoken{సూత్రం}{!!{formula}}.}{}

  \item ప్రతి ప్రతిజ్ఞావాక్య చరం $\Obj p_i$ ఒక (అణు)
  \tetoken{సూత్రం}{!!{formula}}.

  \tagitem{prvNot}{$!A$ \tetoken{ఒక సూత్రం}{!!a{formula}} అయితే, $\lnot !A$
  \tetoken{ఒక సూత్రం}{!!a{formula}}.}{}

  \tagitem{prvAnd}{$!A$, $!B$ \tetoken{సూత్రాలు}{!!{formula}s} అయితే, $(!A \land
  !B)$ \tetoken{ఒక సూత్రం}{!!a{formula}}.}{}

  \tagitem{prvOr}{$!A$, $!B$ \tetoken{సూత్రాలు}{!!{formula}s} అయితే, $(!A \lor !B)$
  \tetoken{ఒక సూత్రం}{!!a{formula}}.}{}

  \tagitem{prvIf}{$!A$, $!B$ \tetoken{సూత్రాలు}{!!{formula}s} అయితే, $(!A \lif !B)$
  \tetoken{ఒక సూత్రం}{!!a{formula}}.}{}

  \tagitem{prvIff}{$!A$, $!B$ \tetoken{సూత్రాలు}{!!{formula}s} అయితే, $(!A \liff !B)$
  \tetoken{ఒక సూత్రం}{!!a{formula}}.}{}

  \item $!A$ \tetoken{ఒక సూత్రం}{!!a{formula}} అయి, $a \in G$ అయితే, $\Knows_a !A$
  \tetoken{ఒక సూత్రం}{!!a{formula}}.

  \item $!A$, $!B$ \tetoken{సూత్రాలు}{!!{formula}s} అయితే, $[!A] !B$
  \tetoken{ఒక సూత్రం}{!!a{formula}}.

  \tagitem{limitClause}{ఇవి తప్ప మరేదీ \tetoken{సూత్రం}{!!a{formula}} కాదు.}{}
\end{enumerate}
\end{defn}

\tetoken{సూత్రం}{!!{formula}} $[!A] !B$కు ఉద్దేశించిన
పఠనం: ``$!A$ను సత్యంగా ప్రకటించిన తరువాత
$!B$ నిలుస్తుంది.'' బహిరంగ ప్రకటనల సందర్భంలో
సామాన్య జ్ఞానం గురించీ మాట్లాడాల్సి రావచ్చు.
కాబట్టి భాషలో సామాన్య జ్ఞాన కారకం~$\CKnows_G
!A$ను కూడా చేర్చవచ్చు.

\end{document}
